%!TEX root = tfg.tex
\chapter{Ejecución de los Planes de Gestión}

\begin{abstract}
Este capítulo describe cómo los planes de gestión definidos previamente se trasladaron a
la ejecución real del proyecto.
\end{abstract}

\section{Ejecución del plan de gestión del alcance}
La ejecución del alcance se articuló sobre FDD, utilizando una visión global del sistema y
una descomposición en funcionalidades manejables. Este enfoque permitió controlar qué
entraba realmente dentro del producto mínimo viable y qué quedaba fuera del alcance del TFG.

\subsection{FDD: Modelo general -- Overall Model}
El modelo general fija la estructura de alto nivel del videojuego. En \textit{Via Inferni}
esto implicó definir la relación entre el núcleo jugable, la generación procedural, la
progresión por círculos, la interfaz y la persistencia de la partida. También se adoptó
una arquitectura modular basada en gestores especializados para mantener una separación
clara de responsabilidades y facilitar la evolución del sistema.

\subsection{FDD: Lista de Funcionalidades -- Feature List}
A partir de ese modelo general se construyó una lista de funcionalidades agrupadas por
dominio: mecánicas del jugador, generación procedural y mundo, entidades y comportamiento,
economía y progresión, e interfaz y \textit{feedback}. Esta descomposición hizo posible
planificar por incrementos verificables y justificar cada avance dentro de la memoria.

\section{Ejecución del plan de gestión del cronograma}
La planificación temporal se ejecutó mediante tres fases principales y una descomposición
interna en iteraciones para la fase de desarrollo. Los hitos se utilizaron como puntos de
control para revisar desviaciones, reajustar la carga de trabajo y decidir si una
funcionalidad debía consolidarse, aplazarse o dividirse en entregas más pequeñas.

\section{Ejecución del plan de gestión de comunicaciones}
La comunicación diaria se apoyó en Discord, mientras que GitHub concentró el seguimiento
formal de incidencias, revisiones y cambios sobre el código. El tablero Kanban de GitHub
Projects se utilizó como herramienta visual de seguimiento, permitiendo distinguir qué
funcionalidades estaban pendientes, en curso, en revisión o cerradas.

\section{Ejecución del plan de gestión de riesgos}
La gestión de riesgos se aplicó de forma preventiva. Los componentes con mayor
incertidumbre, como la generación procedural y la integración entre salas, enemigos e
interfaz, se abordaron tempranamente para detectar problemas antes de que afectaran a una
parte mayor del sistema. Cuando una funcionalidad resultó más costosa de lo previsto, se
subdividió en varias iteraciones para reducir el riesgo de bloqueo.

\section{Ejecución del plan de gestión de la calidad}
La \textit{Definition of Done} se aplicó como criterio de cierre real de cada feature.
Esto evitó considerar terminadas soluciones parciales o aisladas. La validación se realizó
tanto desde el punto de vista técnico como jugable, comprobando que cada incremento
funcionaba en el motor y convivía correctamente con los sistemas ya integrados.

\section{Ejecución del plan de gestión de costes}
La ejecución del plan de costes se reflejó mediante una estimación consolidada del valor
económico del proyecto. El mayor peso recae en el tiempo de trabajo de los dos
desarrolladores, complementado con la amortización del hardware y los costes indirectos del
entorno de desarrollo.

\begin{table}[h!]
    \centering
    \begin{coolTable}{ll}{2}
    {Resumen del presupuesto del proyecto}
    \textbf{Costes de personal (salarios brutos)} & 11.666,66 \euro \\
    \textbf{Costes sociales (30\% sobre bruto)} & 3.500,00 \euro \\
    \textbf{Costes materiales (amortización hardware)} & 125,00 \euro \\
    \textbf{Costes materiales (software y servicios)} & 0,00 \euro \\
    \midrule
    \textbf{Total costes directos} & 15.291,66 \euro \\
    \textbf{Costes indirectos (10\%)} & 1.529,17 \euro \\
    \midrule
    \textbf{TOTAL DEL PROYECTO} & \textbf{16.820,83 \euro} \\
    \end{coolTable}
    \caption{Resumen detallado de la inversión económica total.}
\end{table}
